\subsection*{前言}
\addcontentsline{toc}{subsection}{前言}
本书是对 MIT \texttt{RES.9-009: Introduction to Computational Neuroscience with Neuroblox}（\texttt{January IAP 2025}）的一次 coverage-first 中文重建。全书以 MIT OpenCourseWare 课程页、Neuroblox 官方课程页、MIT Video Productions 的公开录播，以及各单元页面中的官方图、代码与 references 为证据基础，目标不是做一份简短摘要，而是提供一部可自学、可引用、可复习的教材式讲义。

与常见的课程笔记不同，本书把“来源完整性”放在优先级最高的位置。凡是官方页面、字幕、代码块、图示或 reading 中有独立教学价值的内容，都尽量在正文中获得明确落点；无法直接取得的 slide PDF、页面缺图或非教学性片段，则会在工作区的 \texttt{omission\_log.jsonl} 中留下可审计记录，而不是被静默跳过。

本次课程视频的公开 YouTube 入口，最终以 MIT Video Productions 的 \texttt{IAP 2025} playlist 为准，其中包含本课程 11 条编号录播。与此同时，章节顺序与内容边界仍以官方 Neuroblox 课程页的单元结构为主，因为该课程页同时提供了文字讲解、代码与图示，是比纯视频列表更完整的教学证据层。

就“本书是否依照课程视频而来”这一点，可以明确回答：\textbf{是的，但不是机械地只抄视频。} 其中 11 章属于明确的 video-backed 重建，正文直接受官方 YouTube 字幕与页面内容共同约束；另外 3 章（\texttt{08}、\texttt{09}、\texttt{12}）没有官方嵌入视频，因此采用 page-backed 的方式，以官方课程页、图、代码和 reading 为主证据来补足。也就是说，本书整体遵循课程视频与课程页的联合证据，而不是脱离课程材料的自由重写。

从使用角度看，你可以把这本书当成三种东西的结合：
\begin{itemize}
    \item 一部按课程顺序重建的中文教材；
    \item 一套带证据侧车的可审计课程档案；
    \item 一个可以继续增补、修订、重新导出的研究型工作区。
\end{itemize}

如果你需要正式引用本书中的某一部分，建议同时引用对应章节、课程原始单元页以及相关视频或 reading。正文是中文重建与教学化重写，而事实根据与图文出处则应回溯到各章工作区中的 \texttt{source\_manifest.json}、\texttt{figure\_manifest.json} 和 \texttt{eval\_reports/pass\_\#\#.json}。
